\documentclass[11pt]{article}
\usepackage{amsmath,amssymb,amsthm}
\usepackage[T1]{fontenc}
\usepackage{lmodern}
\usepackage[margin=1in]{geometry}
\usepackage{amsmath,amssymb,amsthm}
\usepackage[hidelinks]{hyperref}
\newtheorem{theorem}{Theorem}[section]
\newtheorem{proposition}[theorem]{Proposition}
\newtheorem{lemma}[theorem]{Lemma}
\theoremstyle{remark}
\newtheorem{remark}[theorem]{Remark}
\newcommand{\Z}{\mathbb Z}
\newcommand{\N}{\mathbb N}
\newcommand{\K}{\mathbb K}
\newcommand{\hz}{\mathbf 0}
\newcommand{\hs}{\boxplus}
\DeclareMathOperator{\Diff}{Diff}

\newcommand{\fptitle}{Review of Problem 5 -- MerLean}
\newcommand{\reviewtag}{Reviewer 3}
\usepackage{fp-review}

\begin{document}
\maketitle

\section{Review}

In a ring or a field, the identity $x+y+z = 0$ implies the identity $x^3 + y^3 + z^3 - 3xyz = 0$. This submitted note provides a counterexample for its counterpart for hyperfields, in which the addition is possibly multi-valued. 

The counterexample presented in this note follows the ideas of Hall, Prenowitz, and Connes--Consani. Concretely, Prenowitz and Connes--Consani connected hyperfields with point-line arrangements, and Hall gave constructions of an infinite point-line arrangement from $\mathbb{Z}$. I do not find any mistakes in the proof. However, the historical remarks are only included at the end of Section 3 but not in the main body of the submission. The author cited K. Thas's paper for the infinite point-line arrangement but missed Hall's original paper. If one reads Thas's paper cited in this submission, he or she can quickly realize that the idea presented in this submission originates from Hall. 

All references are in arXiv. The authors should cite the published journal versions. The first reference has been updated by the two authors; the author of this current submission should cite ``Matroids over partial hyperstructures''. 

In terms of individual proofs, the construction of a hyperfield from a projective plane (Propositions 6.1 and 6.2) is almost identical to the construction and proof in Connes--Consani, and the extension construction of the point-line arrangements (Section 5) is almost identical to the construction and proof in Hall. 

For publication in a journal, I recommend a major revision. In particular, the author should remove the proofs that are already written in the literature and only cite the corresponding papers instead. 


\paragraph{Correctness.}
\begin{itemize}
	\item  1: essentially flawless mathematical argument
	\item  0.5: minor errors that are easily corrected
	\item  0: blank solution, or acknowledged partial result
	\item  $-0.5$: major errors
	\item  $-1$: proposed proof makes no progress towards the result.
\end{itemize}
Note that points should be docked for wrong statements/ solutions.

\noindent Score: 1

\paragraph{Exposition.}
\begin{itemize}
	\item  1: well-written and clear solution.
	\item  0: argument that is possible to follow with serious effort
	\item  $-1$: very poor exposition
\end{itemize}

\noindent Score: -1

\paragraph{Attributions.}
\begin{itemize}
	\item  1: solutions contains fully appropriate attributions
		to related work and theorems being used.
	\item  0: moderate issues with missing attributions or hallucinated references
	\item  $-1$: blatant plagiarism or hallucinations
\end{itemize}

\noindent Score: 0

\newpage
\section{Submission with inline comments}


\begin{abstract}
Over a field, a triple with sum zero satisfies the identity
$x^3+y^3+z^3=3xyz$. We show that the corresponding implication for
hyperaddition fails in general, even when the three copies of $-xyz$ are
kept as separate summands. We construct a countably infinite hyperfield
with infinite cyclic multiplicative group in which $(1,g,g^4)$ is a
counterexample. The construction uses translates of a subset of the
integers representing each nonzero integer uniquely as an ordered
difference. We prove the required incidence properties and all
hyperfield axioms directly. We also prove that the implication holds
for quotients of fields, for triples with a singleton pairwise
hypersum, and, under $-1=1$, for triples with equal squares.
\end{abstract}

\section{The question and the result}

A hyperfield has an ordinary multiplication and a multivalued addition:
$x\hs y$ is a nonempty set of elements. A finite list is a
\emph{null sum} if its hypersum contains zero. The question is whether,
in every hyperfield $H$ and for every $x,y,z\in H$,
\begin{equation}\label{eq:question}
 0\in x\hs y\hs z
 \quad\Longrightarrow\quad
 0\in x^3\hs y^3\hs z^3\hs(-xyz)\hs(-xyz)\hs(-xyz).
\end{equation}
Here $-xyz$ denotes the additive inverse of the product $xyz$.
The last three summands are three separate copies of this element;
we do not replace them by a scalar coefficient.

For ordinary fields, the implication follows from
\begin{equation}\label{eq:field-identity}
 x^3+y^3+z^3-3xyz
 =(x+y+z)(x^2+y^2+z^2-xy-yz-zx).
\end{equation}
It also holds in every quotient of a field by a multiplicative subgroup,
because a null triple in the quotient can be lifted to representatives
whose ordinary sum is zero; see Proposition~\ref{prop:quotient}.
The obstacle in a general hyperfield is that membership in a hypersum
does not identify that set with one of its members. Consequently,
ordinary polynomial manipulations cannot be transferred merely by
replacing addition signs with hyperaddition signs.

\begin{theorem}\label{thm:main}
There is a countably infinite hyperfield
\[
 H=\{\hz\}\cup\{g^n:n\in\Z\},\qquad
 g^ag^b=g^{a+b},\qquad g^0=1,
\]
in which every element is its own additive inverse and
\begin{align}
 \hz&\in 1\hs g\hs g^4,\label{eq:witness-premise}\\
 \hz&\notin 1\hs g^3\hs g^{12}\hs g^5\hs g^5\hs g^5.
 \label{eq:witness-failure}
\end{align}
Here $\hz$ is the additive zero, and the $g^n$ are pairwise distinct
nonzero elements. Thus \eqref{eq:question} is false for arbitrary
hyperfields.
\end{theorem}

The geometric method has established precedents. Connes and Consani
relate hyperaddition to projective incidence in
\cite[Proposition~3.1 and Lemma~3.5]{CC}; their
Example~4.7 describes hyperfields from difference sets, including
infinite cyclic examples arising from Hall's construction.
Thas discusses this connection and partial difference sets in
\cite[Sections~2.2--2.3 and 4.1]{Thas}.
We use this method with explicit integer data chosen to make
\eqref{eq:question} fail. The general geometric construction is not a
new assertion here. The proof below supplies the particular difference
set, verifies the hyperfield axioms, and evaluates the six summands.
No classification theorem for hyperfields or projective planes is needed.

Section~\ref{sec:conventions} fixes the hyperfield conventions.
Sections~\ref{sec:differences}--\ref{sec:witness} construct and verify the
counterexample. Section~\ref{sec:positive} gives sufficient conditions
under which the implication does hold.

\section{Hyperaddition conventions}\label{sec:conventions}

We use commutative hyperfields, with the conventions of
\cite[Definitions~2.1, 2.4 and 2.8]{BB}.
For subsets $A,B\subseteq H$, define
\[
 A\hs B=\bigcup_{a\in A,\ b\in B}(a\hs b),
 \qquad tA=\{ta:a\in A\}.
\]
An element used in a subset hypersum stands for its singleton.
A \emph{canonical commutative hypergroup} is a set with a nonempty-valued
operation $\hs$, an element $0$, and the following properties:
\begin{enumerate}
\item $a\hs b=b\hs a$ and $(a\hs b)\hs c=a\hs(b\hs c)$;
\item $0\hs a=\{a\}$;
\item for each $a$, there is a unique element $-a$ with $0\in a\hs(-a)$;
\item $c\in a\hs b$ if and only if $b\in c\hs(-a)$.
\end{enumerate}
\review{The definition of $a \hs b$ should come before the definition of $A \hs B$. }
The last property is called \emph{reversibility}.
A \emph{hyperfield} is such a hypergroup with an abelian group structure
on $H^\times=H\setminus\{0\}$, with identity $1\ne0$, extended by
absorbing zero, such that exact scalar distributivity holds:
\begin{equation}\label{eq:distributivity}
 t(a\hs b)=ta\hs tb\qquad(t,a,b\in H).
\end{equation}
Associativity defines finite hypersums without a choice of parentheses.
Subset hypersums are monotone in each argument by their union definition.

We will use two elementary consequences. Unique inverses give
\begin{equation}\label{eq:zero-test}
 0\in a\hs b\quad\Longleftrightarrow\quad b=-a,
 \qquad
 0\in a\hs b\hs c\quad\Longleftrightarrow\quad -c\in a\hs b.
\end{equation}
For the second equivalence, an intermediate element $u\in a\hs b$
can satisfy $0\in u\hs c$ exactly when $u=-c$.
Distributing $t$ over $0\in a\hs(-a)$ and using uniqueness shows that
$t(-a)=-(ta)$. In particular
\begin{equation}\label{eq:signs}
 -a=(-1)a,\qquad (-1)^2=1,\qquad (-a)^3=-a^3.
\end{equation}
Here $-(-a)=a$ follows already from commutativity and inverse uniqueness.

\section{An integer set with unique differences}\label{sec:differences}

We first construct the incidence structure that will determine
hyperaddition. For a finite set $F\subseteq\Z_{\ge0}$ containing $0$, put
\[
 \Diff(F)=\{d-e:d,e\in F,\ d>e\}.
\]
We say its positive differences are unique if every member of
$\Diff(F)$ has exactly one representation of this form.
The next extension fills a missing difference without changing any
existing representation. Placing both new points far beyond $F$
separates the new cross differences from the old ones.

\begin{lemma}\label{lem:extension}
Suppose $F\subseteq\Z_{\ge0}$ is finite, contains $0$, and has unique
positive differences. Let
\[
 B=\max F,\qquad
 q=\min\bigl(\Z_{>0}\setminus\Diff(F)\bigr),\qquad
 M=2B+q+14.
\]
Then $F'=F\cup\{M,M+q\}$ has unique positive differences, preserves
every old representation, and represents $q$.
\end{lemma}

\begin{proof}
The new positive differences, besides those in $\Diff(F)$, are
\[
 q,\qquad M-d\ (d\in F),\qquad M+q-d\ (d\in F).
\]
Old differences are at most $B$, and $q$ is not an old difference.
Every cross difference is at least
$M-B=B+q+14$, which exceeds both $B$ and $q$.
There are no repetitions within either cross family. An equality
$M-d=M+q-e$ would give $e-d=q$ with $d,e\in F$, contrary to the
choice of $q$. These exhaust the possible collisions. Finally,
$q=(M+q)-M$ is its new representation.
\end{proof}

Start with
\begin{equation}\label{eq:seed}
 D_0=\{0,1,4,6\}.
\end{equation}
Its six positive differences are $1,2,3,4,5,6$, each once.
For $n\ge0$, define
\begin{equation}\label{eq:recursion}
 \begin{split}
 q_n&=\min\bigl(\Z_{>0}\setminus\Diff(D_n)\bigr),\qquad
 B_n=\max D_n,\\
 M_n&=2B_n+q_n+14,\qquad
 D_{n+1}=D_n\cup\{M_n,M_n+q_n\},\qquad
 D=\bigcup_{n\ge0}D_n.
 \end{split}
\end{equation}
The fixed offset $14$ keeps every new point above $13$.
This exclusion will separate the cubic indices from one particular line.

\review{I don't understand what the ``cubic indices'' and the ``one particular line'' mean. }

\begin{lemma}\label{lem:perfect}
The map
\[
 \{(d,e)\in D^2:d\ne e\}\longrightarrow\Z\setminus\{0\},
 \qquad (d,e)\longmapsto d-e
\]
is a bijection. The set $D$ is infinite, contains $0,1,4,6$, and does
not contain $13$.
\end{lemma}

\begin{proof}
Each stage is finite, so $q_n$ exists. Lemma~\ref{lem:extension} applies
inductively. All positive integers below $q_n$ were already represented,
and stage $n+1$ also represents $q_n$. Hence $q_{n+1}>q_n$.
Since $q_0=7$, we have $q_n\ge n+7$; thus $D_n$ represents every
positive integer at most $n+6$. Every positive integer is therefore a
difference in $D$, and reversal covers every negative integer.

If two ordered pairs in $D$ represented the same nonzero difference,
all four entries would belong to a single finite stage. Uniqueness at
that stage rules out distinct representations, also for negative
differences by reversal.

Each extension adds two points above the preceding maximum, so $D$ is
infinite. In the first extension $M_0=33$ and
\[
 D_1=\{0,1,4,6,33,40\}.
\]
All later new points exceed the previous maximum. Since $13\notin D_0$,
we obtain $13\notin D$.
\end{proof}

For $t\in\Z$, call $D+t$ a \emph{line} on the point set $\Z$.
Points are \emph{collinear} if they lie on a common line.
This use of lines is purely an incidence definition.

\begin{lemma}\label{lem:incidence}
Two distinct integer points lie on exactly one line. Two lines with
distinct translation parameters meet in exactly one point and are
different sets. Every line is infinite, and translation by any integer
preserves lines.
\end{lemma}

\begin{proof}
For $a\ne b$, write $a-b=d-e$ with $d,e\in D$, uniquely by
Lemma~\ref{lem:perfect}. The required line is $D+t$, where
$t=a-d=b-e$. Any other line through $a,b$ gives another representation
of $a-b$ and hence the same $d,e,t$.

An intersection point of $D+s$ and $D+t$, for $s\ne t$, has the form
$d+s=e+t$, or equivalently $d-e=t-s$. The latter has exactly one
solution in $D^2$, so the intersection consists of exactly one point.
Since each translate is infinite, the two lines cannot coincide.
Finally, $(D+t)+h=D+(t+h)$.
\end{proof}

Write $L(a,b)$ for the unique line through distinct integer points
$a,b$. In particular,
\begin{equation}\label{eq:translation}
 L(a+h,b+h)=L(a,b)+h.
\end{equation}
Only the properties proved in Lemma~\ref{lem:incidence}, rather than an
external incidence theorem, will be used below.

\section{The hyperfield and its axioms}\label{sec:hyperfield}

Let
\[
 H=\{\hz\}\cup\{g^n:n\in\Z\},
\]
where all the displayed symbols are distinct. The additive zero $\hz$
is different from the integer index $0$ and from $g^0=1$.
Transport each integer line to $H^\times$ by $n\mapsto g^n$.
For distinct $A=g^a,B=g^b$, write
\[
 \ell(A,B)=\{g^c:c\in L(a,b)\}.
\]
The incidence properties of Lemma~\ref{lem:incidence} hold for these
transported lines as well.

Define hyperaddition by
\begin{equation}\label{eq:addition}
 \hz\hs u=u\hs\hz=\{u\},\qquad
 A\hs A=\{\hz,A\},\qquad
 A\hs B=\ell(A,B)\setminus\{A,B\}\quad(A\ne B),
\end{equation}
where $A,B$ in the last two rules are nonzero. This is the usual
line-based addition from \cite[Proposition~3.1]{CC}, here applied to the
specific integer lines. The removal of the two endpoints makes three
distinct nonzero elements add to zero precisely when they are collinear.
We verify the laws directly, including the associativity calculation
needed to evaluate the counterexample.

\begin{proposition}\label{prop:hypergroup}
The operation \eqref{eq:addition} makes $H$ a canonical commutative
hypergroup, with $-u=u$ for all $u\in H$. For noncollinear nonzero points
$A,B,C$, its three-term hypersum is
\begin{equation}\label{eq:triangle}
 A\hs B\hs C
 =H^\times\setminus\bigl(\ell(A,B)\cup\ell(A,C)\cup\ell(B,C)\bigr).
\end{equation}
\end{proposition}

\begin{proof}
Every line is infinite, so every addition value is nonempty.
Commutativity and the zero identity are immediate. The rules give
\begin{equation}\label{eq:self-inverse}
 \hz\in u\hs v\quad\Longleftrightarrow\quad u=v.
\end{equation}
Thus each element is its own unique additive inverse.

To verify reversibility, examine the relation $w\in u\hs v$.
If one entry is zero, the permitted triples are the permutations of
$(\hz,u,u)$, including the all-zero case. For three nonzero entries,
the permitted patterns are that all three are equal, or that they are
distinct and collinear. Each condition is invariant under permutation.
Since negation is the identity, this symmetry is exactly reversibility.

It remains to prove associativity. We compute the left-associated
hypersum for each possible configuration; the resulting expressions
also determine the other parenthesization.
A zero entry reduces the assertion to the identity rule. Suppose
therefore that $A,B,C$ are nonzero.

If all three entries equal $A$, both parenthesizations give
$\{\hz,A\}$. If exactly two entries agree, denote the repeated point
by $A$ and the other by $B$, and put $\ell=\ell(A,B)$. Then
\[
 (A\hs A)\hs B
 =\{B\}\cup\bigl(\ell\setminus\{A,B\}\bigr)
 =\ell\setminus\{A\},
\]
whereas
\[
 A\hs(A\hs B)
 =\bigcup_{U\in\ell\setminus\{A,B\}}
       \bigl(\ell\setminus\{A,U\}\bigr)
 =\ell\setminus\{A\}.
\]
For the last equality, any point of $\ell$ other than $A$ survives by
choosing $U$ different from that point; the infinite line supplies such
choices. Commutativity covers the other placements of $B$.

Suppose next that $A,B,C$ are distinct and lie on a line $\ell$.
In $(A\hs B)\hs C$, the intermediate choice $U=C$ contributes
$\{\hz,C\}$, while each $U\in\ell\setminus\{A,B,C\}$ contributes
$\ell\setminus\{U,C\}$. There are at least two such $U$, so
\[
 (A\hs B)\hs C=\ell\cup\{\hz\}.
\]
The same calculation applies with the entries permuted and proves
associativity in this case.

Finally let $A,B,C$ be noncollinear. We prove
\eqref{eq:triangle} first for $(A\hs B)\hs C$.
An intermediate point $U$ belongs to
$\ell(A,B)\setminus\{A,B\}$, so $U\ne C$. A point
$V\in U\hs C$ lies on $\ell(U,C)$ and differs from $U,C$.
The lines $\ell(U,C)$ and $\ell(A,B)$ are distinct and meet only at
$U$. Moreover, $\ell(U,C)$ differs from both $\ell(A,C)$ and
$\ell(B,C)$: equality with either would give two intersections with
$\ell(A,B)$. Its intersection with each of these latter lines is
therefore only $C$. Hence $V$ lies outside all three lines on the
right of \eqref{eq:triangle}. Also $\hz$ cannot occur, since $U\ne C$.

Conversely, let $V$ be outside those three lines. The line $\ell(C,V)$
meets $\ell(A,B)$ at a unique point $U$. This point is neither $A$
nor $B$, since otherwise $V$ would belong to $\ell(A,C)$ or
$\ell(B,C)$. Also $V\ne U,C$. Thus
$U\in A\hs B$ and $V\in U\hs C$, proving the reverse inclusion.
The right side of \eqref{eq:triangle} is symmetric in $A,B,C$.
Applying the same computation to $(B\hs C)\hs A$ and using
commutativity yields $A\hs(B\hs C)$, as required.
\end{proof}

\begin{proposition}\label{prop:hyperfield}
Define $g^ag^b=g^{a+b}$ and make $\hz$ absorbing. Together with
\eqref{eq:addition}, these operations make $H$ a countably infinite
hyperfield.
\end{proposition}

\begin{proof}
The nonzero elements form an infinite cyclic abelian group, with
identity $g^0=1$ and inverse $(g^a)^{-1}=g^{-a}$. In particular,
$1\ne\hz$. Multiplication by $g^h$ translates all integer indices by
$h$. Equation~\eqref{eq:translation} shows that it carries the line
through $g^a,g^b$ to the line through $g^{a+h},g^{b+h}$ and carries
the removed endpoints to the removed endpoints. It therefore preserves
the distinct nonzero case of \eqref{eq:addition} exactly. The equal
and zero cases are preserved by inspection of their formulas. This
proves \eqref{eq:distributivity} for every nonzero scalar. For scalar
$\hz$, both sides are $\{\hz\}$, using nonemptiness of addition values.
Proposition~\ref{prop:hypergroup} supplies all remaining axioms.
The underlying set is countably infinite.
\end{proof}

\section{Evaluation of the counterexample}\label{sec:witness}

\begin{proof}[Proof of Theorem~\ref{thm:main}]
Take the hyperfield of Proposition~\ref{prop:hyperfield}.
The indices $0,1,4$ are distinct members of $D$. Thus
$g^4\in 1\hs g$, and $\hz\in g^4\hs g^4$ proves
\eqref{eq:witness-premise}.

For $(x,y,z)=(1,g,g^4)$, the cubes are $1,g^3,g^{12}$.
Every element is its own negative, so each mixed summand is
$-xyz=xyz=g^5$. To locate their hypersum, observe that
\begin{equation}\label{eq:key-line}
 L(0,3)=D-1,\qquad 5\in D-1,\qquad 12\notin D-1.
\end{equation}
Indeed, $1,4\in D$ identify the line through $0,3$;
$6\in D$ gives the second assertion; and $13\notin D$ gives the third.
Consequently $1,g^3,g^{12}$ are noncollinear. By
\eqref{eq:triangle}, the set
\[
 T=1\hs g^3\hs g^{12}
\]
contains neither $\hz$ nor $g^5$, since $g^5$ lies on the line
through $1,g^3$.

On the other hand, the equal-element addition rule gives
\[
 g^5\hs g^5\hs g^5
 =\{\hz,g^5\}\hs g^5=\{\hz,g^5\}.
\]
Associativity identifies the required six-term hypersum with
$T\hs\{\hz,g^5\}$. By \eqref{eq:self-inverse},
\[
 \hz\in T\hs\{\hz,g^5\}
 \quad\Longleftrightarrow\quad
 T\cap\{\hz,g^5\}\ne\varnothing.
\]
The intersection is empty, proving \eqref{eq:witness-failure}.
\end{proof}

\section{Cases in which the implication holds}\label{sec:positive}

The failure above does not apply to every hyperfield or every null
triple. We record several sufficient conditions and prove them using
only the stated hyperfield laws. In this section $H$ is again an
arbitrary hyperfield and its zero is denoted by $0$.

\begin{lemma}\label{lem:normalization}
The implication \eqref{eq:question} holds if at least one coordinate
is zero. If $xyz\ne0$, set
\[
 t=-z^{-1},\qquad a=tx,\qquad b=ty.
\]
Then its premise is equivalent to $1\in a\hs b$, and its conclusion
is equivalent to
\begin{equation}\label{eq:normalized}
 1\in a^3\hs b^3\hs ab\hs ab\hs ab.
\end{equation}
Conversely, a failure of \eqref{eq:normalized} with $1\in a\hs b$
yields a counterexample $(a,b,-1)$ to \eqref{eq:question}.
\end{lemma}

\begin{proof}
If $z=0$, then $y=-x$ by \eqref{eq:zero-test}; the mixed summands
vanish and the cubes form an opposite pair by \eqref{eq:signs}.
Symmetry covers all positions of a zero, including the all-zero triple.

For nonzero coordinates, multiplication by $t$ is a bijection fixing
zero and preserving hypersums, and $tz=-1$. The premise is therefore
equivalent to $0\in a\hs b\hs(-1)$, or $1\in a\hs b$.
Multiplying the six-term conclusion by $t^3$ gives precisely
\[
 0\in a^3\hs b^3\hs(-1)\hs ab\hs ab\hs ab,
\]
because $(tz)^3=-1$ and $t^3(-xyz)=-(tx)(ty)(tz)=ab$.
The zero test converts this to \eqref{eq:normalized}. Every step is
reversible. Substituting $(a,b,-1)$ proves the final assertion.
\end{proof}

\begin{proposition}\label{prop:quotient}
Let $F$ be a commutative field and $G\le F^\times$ a subgroup.
Then \eqref{eq:question} holds in the quotient hyperfield $F/G$,
in every characteristic.
\end{proposition}

\begin{proof}
Write $[r]=rG$ for $r\ne0$, and $[0]=\{0\}$. Multiplication is
$[r][s]=[rs]$, and hyperaddition is
\begin{equation}\label{eq:quotient}
 [r]\hs[s]=\{[r'+s']:r'\in[r],\ s'\in[s]\}.
\end{equation}
This is the quotient construction of
\cite[Proposition~2.5]{CC} and \cite[Remark~2.7]{BB}.
These results apply because $F$ is a commutative ring and $G$ is a
subgroup of its units; every nonzero class is invertible since $F$
is a field. The additive inverse is $-[r]=[-r]$, as
$[0]\in[r]\hs[-r]$ and inverses are unique.

Suppose $0\in x\hs y\hs z$. By \eqref{eq:zero-test},
$-z\in x\hs y$, so \eqref{eq:quotient} provides representatives
$r\in x$, $s\in y$ with $[r+s]=-z$. Set $q=-(r+s)$.
Then $q\in z$ and $r+s+q=0$. Identity~\eqref{eq:field-identity} gives
\[
 r^3+s^3+q^3-rsq-rsq-rsq=0.
\]
For any finite list in $F$, the class of its ordinary sum belongs to
the hypersum of the classes: the two-term statement is
\eqref{eq:quotient}, and induction appends one term at a time.
Projecting the displayed equality proves \eqref{eq:question}.
No division by $3$ is involved.
\end{proof}

\begin{proposition}\label{prop:singleton}
Suppose $0\in x\hs y\hs z$ in a hyperfield. If any of
$x\hs y$, $y\hs z$, $z\hs x$ is a singleton, then the conclusion
of \eqref{eq:question} holds.
\end{proposition}

\begin{proof}
Lemma~\ref{lem:normalization} treats zero coordinates.
Otherwise permute the singleton pair to $x,y$. By
\eqref{eq:zero-test} its sum is $\{-z\}$, so the same normalization
reduces the claim to $a\hs b=\{1\}$.
Scaling $1\in a\hs b$ by $a$ and $b$ gives
\[
 a\in a^2\hs ab,\qquad b\in ab\hs b^2.
\]
Monotonicity and associativity yield
$1\in a^2\hs b^2\hs ab\hs ab$.
Scaling instead by $a^2$ and $b^2$ gives
\[
 a^2\in a^3\hs a^2b,\qquad b^2\in ab^2\hs b^3.
\]
Substituting these memberships into the preceding hypersum yields
\[
 1\in a^3\hs b^3\hs ab\hs ab\hs(a^2b\hs ab^2).
\]
Here exact distributivity and the singleton hypothesis give the
set equality $a^2b\hs ab^2=ab(a\hs b)=\{ab\}$.
This proves \eqref{eq:normalized}, and scaling back proves the claim.
\end{proof}

\begin{proposition}\label{prop:squares}
Suppose $-1=1$, $x^2=y^2=z^2$, and $0\in x\hs y\hs z$.
Then the conclusion of \eqref{eq:question} holds.
\end{proposition}

\begin{proof}
The zero-coordinate case follows from Lemma~\ref{lem:normalization}.
Otherwise put $k=x^2=y^2=z^2\ne0$, $p=xyz$, and $q=p/k$.
Negation is the identity by \eqref{eq:signs}. Thus
$z\in x\hs y$, and scaling this membership by $xy/k$ gives
\[
 q\in y\hs x,
\]
since $(xy/k)x=y$ and $(xy/k)y=x$.
Cyclically, $q\in y\hs z$ and $q\in z\hs x$.
Reversibility now gives
\[
 y\in x\hs q,\qquad z\in y\hs q,\qquad x\in z\hs q.
\]
Consequently, by monotonicity,
\[
 0\in y\hs z\hs x
 \subseteq (x\hs q)\hs(y\hs q)\hs(z\hs q).
\]
Multiply by $k$. Since $kx=x^3$, $ky=y^3$, $kz=z^3$, and
$kq=p=-xyz$, the resulting membership is the required conclusion.
\end{proof}

\begin{remark}\label{rem:cancellation}
The singleton equality used in Proposition~\ref{prop:singleton} is
essential to that argument. Knowing only $ab\in a^2b\hs ab^2$
does not allow a membership in the larger hypersum to be retained
after replacing this set by $\{ab\}$. Cancellation of opposite
summands is likewise unavailable in general. For example, in the
Krasner hyperfield $\K=\{0,1\}$ with $1\hs1=\{0,1\}$,
\[
 0\in1\hs1\hs1=\{0,1\},\qquad 0\notin\{1\}.
\]
Here $-1=1$, so deleting an opposite pair destroys zero membership.
This familiar hyperfield is described in \cite[Example~2.10]{BB};
it also equals $\mathbb F_3/\mathbb F_3^\times$ under
\eqref{eq:quotient}. Thus the cancellation example itself satisfies
the cubic implication by Proposition~\ref{prop:quotient}.
\end{remark}

Any counterexample therefore has three nonzero coordinates and no
singleton pairwise hypersum. The witness in Theorem~\ref{thm:main}
has infinite pairwise hypersums and distinct squares. The construction
is infinite and does not settle whether a finite counterexample exists.
The quotient result also shows that this particular hyperfield cannot
be isomorphic to a quotient $F/G$ of a field: such an isomorphism would
preserve the premise and all six terms of the failed conclusion.

\begin{thebibliography}{9}
\bibitem{BB}
M. Baker and N. Bowler,
\emph{Matroids over hyperfields},
arXiv:1601.01204v5 (2017),
\url{https://arxiv.org/abs/1601.01204v5}.

\bibitem{CC}
A. Connes and C. Consani,
\emph{The hyperring of ad\`ele classes},
arXiv:1001.4260v2 (2010),
\url{https://arxiv.org/abs/1001.4260v2}.

\bibitem{Thas}
K. Thas,
\emph{Hyperfield extensions, characteristic one and the Connes--Consani
plane connection},
arXiv:1407.0607v1 (2014),
\url{https://arxiv.org/abs/1407.0607v1}.
\end{thebibliography}

\end{document}
